\section{Introduction}\label{sec:intro}

\IEEEPARstart{A} user who asks a model provider to delete their data receives, at best, an
assurance that an unlearning routine was run. Regulators and users increasingly ask for
evidence instead, and the evidence that can be produced without trusting the provider's
own report is statistical: plant records whose fate the auditor knows, ask for their
deletion, and test whether the deployed model still favours them. The secret-sharer
methodology \citep{carlini2019secret}, privacy auditing with canaries
\citep{steinke2023privacy}, copyright traps \citep{meeus2024traps} and membership-inference
tests on a forget set \citep{shi2024muse,hayes2024inexact} all follow this pattern. The most
recent instance, VOUCH \citep{dang2026vouch}, plants \emph{paired ghost canaries}, two
exchangeable records of which a fair coin admits one into training and later into the
forget set. Under exact unlearning the sign of the score difference between the twins is
then a fair coin, and betting e-processes turn the signs into an anytime-valid certificate
with an exact null.

Every audit of this kind certifies a statement about the planted records. The statement
that matters to users is about their own records. The two coincide only if the provider
processes canaries exactly as it processes everything else, and audits typically assume
it; VOUCH calls the assumption \emph{honest but verifiable}. A provider with an incentive to
pass the audit cheaply, by unlearning carefully what the auditor will check and cutting
corners elsewhere, violates it. This paper asks how much such a provider can hide, what
it needs in order to hide it, and what an auditor can do about it.

\paragraph{Threat.} We consider a \emph{canary-aware} provider: it sees the training corpus
and the forget set, including the planted records, and may run any rule on their content to
decide how thoroughly to unlearn each record. It does not see the auditor's coins. The
provider passes the audit if the canary advantage stays below the auditor's tolerance, and
it gains by leaving as much organic data as possible un-unlearned. Provenance attacks, in
which canaries reach the provider through a distinguishable channel, and providers that
deviate from the committed pipeline altogether are out of scope; the former must be closed
by the injection protocol and the latter by attestation \citep{eisenhofer2025verifiable}.

\paragraph{Contributions.} The contributions are a bound, its measurement on two canary
designs, and a canary design meant to remove the detectability term, together with what it costs.

\begin{itemize}
\item \textbf{An evasion bound and its sharp form.} For audits whose statistic averages a
bounded per-pair leakage over randomly included canaries, the gap between the canary advantage
and the advantage the same pipeline run would show on organic records planted in the same way
is at most twice the distinguishing advantage between canaries and organic records
(Theorem~\ref{thm:detect}). For a provider whose leakage is one-sided and determined by the
trained record, the bound is sharp: the largest organic advantage compatible with a canary
advantage of at most $\eps$ is the Neyman--Pearson power of testing canary records against
organic records at level $\eps$, which is $\eps$ when the two coincide and $1$ when they are
mutually singular (Theorem~\ref{thm:sharp}). Every canary detector therefore yields a provider
strategy, and every strategy yields a detector.

\item \textbf{The canary designs we tested are fully evadable.} A character $n$-gram classifier
recognises VOUCH's random-string canaries and semantic canaries built from their own
templates with advantage one, because of artefacts of their generators. In a simulation
driven by its decisions, a selective provider holds its canary miss rate, an upper bound on
the canary advantage, at the audit's tolerance of $0.2$ while skipping at least $99\%$ of
organic records (Section~\ref{sec:exp-detect}).

\item \textbf{Organic-matched canaries with an exact null.} Paired canaries stating two
incompatible facts keep the audit's exact null for every score that compares the facts
under a common prompt, including scores on paraphrased questions never trained
(Proposition~\ref{prop:semantic}). Drawing both twins from the organic record generator
makes each twin distributed as an organic record, and our detectors do not separate such
canaries from organic records beyond their own noise.

\item \textbf{The detectability--power trade-off.} The bound implies that the canary advantage
can exceed the counterfactual organic advantage only by twice the detectability
(Corollary~\ref{cor:tradeoff}). On Pythia-410M at our training budget, the audit with matched
canaries seen once is valid but has no power, and it certifies a model that never unlearned in
$97\%$ of seed-by-order runs (Section~\ref{sec:exp-matched}); a canary-aware provider would not
need to evade it.

\item \textbf{Membership-inference audits and incomplete forget sets.} A membership-inference
audit that compares canaries with an organic holdout is decided by the distribution shift
between the two: its AUC on a retrained, exactly unlearned model is $0.000$, so a two-sided
test accuses that model in every seed. A paired test on the same records is valid and about as
powerful as VOUCH (Section~\ref{sec:exp-baselines}). When the forget set omits a form in which a
fact was trained, adding paraphrase scores to the audit revokes up to $49\%$ of the runs that a
literal audit certifies in that cell (Section~\ref{sec:exp-squeeze}).
\end{itemize}

\paragraph{Novelty relative to the closest work.} The closest prior work is VOUCH
\citep{dang2026vouch}, which introduced paired ghost canaries and the anytime-valid
certificate that we use as the running example. VOUCH proves that recognising canaries does
not invalidate the certificate on the cohort, and its threat model assumes that the
provider does not special-case canaries; it measures that its canaries are easy to
recognise but does not ask what recognition buys. This paper drops that assumption. The
bound of Theorems~\ref{thm:detect} and~\ref{thm:sharp}, the selective-provider evaluation,
the semantic and organic-matched canary designs with their exact null, the
detectability--power trade-off and the comparison with membership-inference audits are new.
\citet{zhang2024verification} show that a dishonest provider can defeat existing
verification strategies; we quantify the defeat for planted-canary audits as a function of
canary detectability, and show where it can and cannot be closed.
\citet{duan2024membership} and \citet{das2024blind} observe that membership-inference
benchmarks with shifted non-members measure the shift; Theorem~\ref{thm:detect} is the
certificate-level form of that observation, and Section~\ref{sec:exp-baselines} shows its
consequence for unlearning audits: false accusation of an honest provider.

\paragraph{Scope.} Data can bound the distinguishing advantage only from below, so the
statement about organic records is conditional on the strongest detector anyone can build;
we report estimates for the detectors we built. The organic advantage is a counterfactual
quantity: the leakage of organic records planted like canaries, not of the particular records
users deleted. All records are fictitious biographies from one domain,
and the models have at most one billion parameters. The selective provider is simulated
from detector decisions rather than run end to end. Section~\ref{sec:discussion} discusses
these limits.

\paragraph{Organisation.} Section~\ref{sec:model} defines the audit and the threat model.
Section~\ref{sec:bound} proves the bound. Section~\ref{sec:design} presents canary designs
and the cost of continuous monitoring. Section~\ref{sec:experiments} reports the
experiments, Section~\ref{sec:related} discusses related work and
Section~\ref{sec:discussion} the limitations. Proofs not given in the text and additional
results are in the supplementary material.
